\inicializarSeccion{Modificaci\'on del c\'odigo}

\begin{lstlisting}[style=Matlab-editor]
%% Parametros que puedes cambiar
q   = 1e-9;    % Carga electrica de cada punto [C]
d   = 0.10;    % Separacion entre lineas de carga [m]
Np  = 8;       % Numero de cargas positivas
Nn  = 8;       % Numero de cargas negativas
Lp  = 0.20;    % Longitud de la linea de cargas positivas [m]
Ln  = 0.20;    % Longitud de la linea de cargas negativas [m]
N   = 30;      % Puntos en la malla
Lim = 0.30;    % Tamano del espacio simulado [m]

%% Constante de Coulomb
k = 8.99e9;

%% Creamos la malla de puntos en 2D
x = linspace(-Lim, Lim, N);
y = linspace(-Lim, Lim, N);
[X, Y] = meshgrid(x, y);

%% Posicion de cada carga positiva (linea horizontal en y = d/2)
y_pos = d/2;
x_pos = zeros(1, Np);
for i = 1:Np
    x_pos(i) = -Lp/2 + (i-1) * Lp/(Np-1);
end

%% Posicion de cada carga negativa (linea horizontal en y = -d/2)
y_neg = -d/2;
x_neg = zeros(1, Nn);
for j = 1:Nn
    x_neg(j) = -Ln/2 + (j-1) * Ln/(Nn-1);
end

%% Campo total inicializado en cero
Ex = zeros(size(X));
Ey = zeros(size(X));

%% Sumamos el campo electrico de cada carga positiva
for i = 1:Np
    r_pos = sqrt((X - x_pos(i)).^2 + (Y - y_pos).^2);
    r_pos(r_pos < 1e-10) = 1e-10;
    Ex = Ex + k * q * (X - x_pos(i)) ./ r_pos.^3;
    Ey = Ey + k * q * (Y - y_pos) ./ r_pos.^3;
end

%% Sumamos el campo electrico de cada carga negativa
for j = 1:Nn
    r_neg = sqrt((X - x_neg(j)).^2 + (Y - y_neg).^2);
    r_neg(r_neg < 1e-10) = 1e-10;
    Ex = Ex + k * (-q) * (X - x_neg(j)) ./ r_neg.^3;
    Ey = Ey + k * (-q) * (Y - y_neg) ./ r_neg.^3;
end

%% Normalizamos para que todas las flechas tengan el mismo tamano
E_mag = sqrt(Ex.^2 + Ey.^2);
E_mag(E_mag == 0) = 1;
Ex_norm = Ex ./ E_mag;
Ey_norm = Ey ./ E_mag;

%% Graficamos el campo electrico
figure;
quiver(X, Y, Ex_norm, Ey_norm, 0.5, 'b', 'LineWidth', 1.2);
hold on;

%% Marcamos la posicion de cada carga positiva
for i = 1:Np
    plot(x_pos(i), y_pos, 'ro', 'MarkerSize', 8, 'MarkerFaceColor', 'r');
end
text(0, y_pos + 0.01, '+q', 'FontSize', 12, 'Color', 'r', 'FontWeight', 'bold');

%% Marcamos la posicion de cada carga negativa
for j = 1:Nn
    plot(x_neg(j), y_neg, 'bo', 'MarkerSize', 8, 'MarkerFaceColor', 'b');
end
text(0, y_neg + 0.01, '-q', 'FontSize', 12, 'Color', 'b', 'FontWeight', 'bold');

%% Formato de la figura
title(sprintf('Campo electrico | Np = %d, Nn = %d, d = %.2f m', Np, Nn, d), 'FontSize', 13);
xlabel('x [m]', 'FontSize', 12);
ylabel('y [m]', 'FontSize', 12);
axis equal;
axis([-Lim Lim -Lim Lim]);
grid on;
hField = findobj(gca, 'Type', 'Quiver'); % campo eléctrico (flechas)
hPos   = findobj(gca, 'Type', 'Line', '-and', 'Marker', 'o', 'MarkerFaceColor', 'r'); % cargas +q
hNeg   = findobj(gca, 'Type', 'Line', '-and', 'Marker', 'o', 'MarkerFaceColor', 'b'); % cargas -q
legend([hField(1), hPos(1), hNeg(1)], {'Campo eléctrico', 'Carga +q', 'Carga -q'}, 'Location', 'northeast');

hold off;
\end{lstlisting}